% Seminar 2 - Page 4. Fragment for % Seminar 2 - Page 4. Fragment for % Seminar 2 - Page 4. Fragment for % Seminar 2 - Page 4. Fragment for \input{seminars/pages/seminar2-pagina4-en}.
% No preamble and no forced page break.
% EDITORIAL NOTES:
% - This page continues the work begun on Page 3: it first returns to
%   parts (a) and (b), then begins the proof of (c).
% - The equality in (b), for injective f, tacitly assumes a non-empty
%   index family; the manuscript does not discuss the empty case.
% - In the finite example A={-1,0,1}, f(x)=x^2, the handwriting seems
%   to say f(X_1)=A, although f({-1,1})={1}. The apparent A is retained.
% - Right-margin interval calculations seem to concern f(x)=x^2, but
%   the domain is not restated. The manuscript appears to say
%   f([-1,3])=[1,9]; mathematically its image would be [0,9].
%   That apparent source expression is retained and flagged here.
% - Interval brackets and certain endpoints in the boxed calculation
%   are provisional readings of the handwriting.
% - The final proof (c) ends mid-argument; no continuation or part (d)
%   has been added.

\subsection*{Seminar 2 - Page 4}

\noindent\underline{Solution (continued)}

\medskip
\noindent $a)$ \underline{$\subseteq$}.
Let $y\in f\left(\bigcup_{i\in I}X_i\right)$. Then
\[
\exists\,x\in\bigcup_{i\in I}X_i\quad\text{such that}\quad f(x)=y.
\]
\noindent Thus there is $i_0\in I$ with $x\in X_{i_0}$ and $f(x)=y$, hence
\[
 y\in f(X_{i_0})\subseteq\bigcup_{i\in I}f(X_i).
\]
\noindent \underline{$\supseteq$}.
There is $i_0\in I$ such that
$y\in f(X_{i_0})$, whence
\[
 y\in f\left(\bigcup_{i\in I}X_i\right).
\]

\medskip
\noindent $b)$ Let
\[
 y\in f\left(\bigcap_{i\in I}X_i\right).
\]
\noindent There is $x\in\bigcap_{i\in I}X_i$ with $f(x)=y$, so
\[
 y\in\bigcap_{i\in I}f(X_i).
\]
\noindent Suppose $f$ is injective. If
$y\in\bigcap_{i\in I}f(X_i)$, then
$y\in f(X_i)$ for every $i\in I$. Thus
\[
 (\forall i\in I)(\exists a_i\in X_i):\ f(a_i)=y.
\]
\noindent By injectivity, $a_i=a_j$ for all $i,j\in I$.
Denote this common element by $a$. Then
\[
 f(a)=y,\qquad a\in\bigcap_{i\in I}X_i,
 \qquad y\in f\left(\bigcap_{i\in I}X_i\right).
\]

\medskip
\noindent\underline{Example.} The manuscript also gives
\[
 A=\{-1,0,1\},\qquad B=\{0,1\},\qquad f(x)=x^2,
\]
\[
 X_1=\{-1,1\},\qquad X_2=\{-1,0\},
\]
\[
 f(X_1)=A,\qquad f(X_2)=\{0,1\}.
\]

% Small sketch of the parabola y=x^2, after the marginal drawing.
\begin{center}
\begin{tikzpicture}[>=stealth,scale=.77]
  \draw[->] (-1.8,0)--(3.0,0) node[right] {$x$};
  \draw[->] (0,-.25)--(0,3.45) node[above] {$y$};
  \foreach \x in {-1,0,1,2,3}{\draw (\x,.06)--(\x,-.06) node[below] {\scriptsize $\x$};}
  \draw[domain=-1.35:2.35,samples=70,smooth,thick] plot (\x,{.5*\x*\x});
\end{tikzpicture}
\end{center}

% Right-margin examples: interval brackets and some endpoints are uncertain.
\noindent\underline{Marginal calculations (provisional readings):}
\[
\begin{aligned}
 f([1,2])&=[1,4], & f([-1,3])&=[1,9],\\
 f([1,2])&=f([1,3]\cap[1,2])=[1,4].
\end{aligned}
\]
\[
\boxed{\begin{aligned}
 f([-1,0])&=[0,1],\\
 f([0,2])&=[0,4],\\
 f([-1,0]\cap[0,2])&=f(0)=0,\\
 f([-1,0])\cap f([0,2])&=[0,1].
\end{aligned}}
\]

\medskip
\noindent $c)$ \underline{$\subseteq$}.
Let
\[
 x\in f^{-1}\left(\bigcup_{j\in J}Y_j\right).
\]
\noindent Then
\[
 f(x)\in\bigcup_{j\in J}Y_j
 \quad\Longrightarrow\quad
 (\exists j_0\in J):\ f(x)\in Y_{j_0}
 \quad\Longrightarrow\quad
 x\in f^{-1}(Y_{j_0})
 \subseteq\bigcup_{j\in J}f^{-1}(Y_j).
\]
\noindent \underline{$\supseteq$}.
If $x\in\bigcup_{j\in J}f^{-1}(Y_j)$, there is $j_0\in J$
with $x\in f^{-1}(Y_{j_0})$, and therefore
\[
 f(x)\in Y_{j_0}\ \Longrightarrow\ f(x)\in\bigcup_{j\in J}Y_j.
\]
% The handwritten page stops here.
.
% No preamble and no forced page break.
% EDITORIAL NOTES:
% - This page continues the work begun on Page 3: it first returns to
%   parts (a) and (b), then begins the proof of (c).
% - The equality in (b), for injective f, tacitly assumes a non-empty
%   index family; the manuscript does not discuss the empty case.
% - In the finite example A={-1,0,1}, f(x)=x^2, the handwriting seems
%   to say f(X_1)=A, although f({-1,1})={1}. The apparent A is retained.
% - Right-margin interval calculations seem to concern f(x)=x^2, but
%   the domain is not restated. The manuscript appears to say
%   f([-1,3])=[1,9]; mathematically its image would be [0,9].
%   That apparent source expression is retained and flagged here.
% - Interval brackets and certain endpoints in the boxed calculation
%   are provisional readings of the handwriting.
% - The final proof (c) ends mid-argument; no continuation or part (d)
%   has been added.

\subsection*{Seminar 2 - Page 4}

\noindent\underline{Solution (continued)}

\medskip
\noindent $a)$ \underline{$\subseteq$}.
Let $y\in f\left(\bigcup_{i\in I}X_i\right)$. Then
\[
\exists\,x\in\bigcup_{i\in I}X_i\quad\text{such that}\quad f(x)=y.
\]
\noindent Thus there is $i_0\in I$ with $x\in X_{i_0}$ and $f(x)=y$, hence
\[
 y\in f(X_{i_0})\subseteq\bigcup_{i\in I}f(X_i).
\]
\noindent \underline{$\supseteq$}.
There is $i_0\in I$ such that
$y\in f(X_{i_0})$, whence
\[
 y\in f\left(\bigcup_{i\in I}X_i\right).
\]

\medskip
\noindent $b)$ Let
\[
 y\in f\left(\bigcap_{i\in I}X_i\right).
\]
\noindent There is $x\in\bigcap_{i\in I}X_i$ with $f(x)=y$, so
\[
 y\in\bigcap_{i\in I}f(X_i).
\]
\noindent Suppose $f$ is injective. If
$y\in\bigcap_{i\in I}f(X_i)$, then
$y\in f(X_i)$ for every $i\in I$. Thus
\[
 (\forall i\in I)(\exists a_i\in X_i):\ f(a_i)=y.
\]
\noindent By injectivity, $a_i=a_j$ for all $i,j\in I$.
Denote this common element by $a$. Then
\[
 f(a)=y,\qquad a\in\bigcap_{i\in I}X_i,
 \qquad y\in f\left(\bigcap_{i\in I}X_i\right).
\]

\medskip
\noindent\underline{Example.} The manuscript also gives
\[
 A=\{-1,0,1\},\qquad B=\{0,1\},\qquad f(x)=x^2,
\]
\[
 X_1=\{-1,1\},\qquad X_2=\{-1,0\},
\]
\[
 f(X_1)=A,\qquad f(X_2)=\{0,1\}.
\]

% Small sketch of the parabola y=x^2, after the marginal drawing.
\begin{center}
\begin{tikzpicture}[>=stealth,scale=.77]
  \draw[->] (-1.8,0)--(3.0,0) node[right] {$x$};
  \draw[->] (0,-.25)--(0,3.45) node[above] {$y$};
  \foreach \x in {-1,0,1,2,3}{\draw (\x,.06)--(\x,-.06) node[below] {\scriptsize $\x$};}
  \draw[domain=-1.35:2.35,samples=70,smooth,thick] plot (\x,{.5*\x*\x});
\end{tikzpicture}
\end{center}

% Right-margin examples: interval brackets and some endpoints are uncertain.
\noindent\underline{Marginal calculations (provisional readings):}
\[
\begin{aligned}
 f([1,2])&=[1,4], & f([-1,3])&=[1,9],\\
 f([1,2])&=f([1,3]\cap[1,2])=[1,4].
\end{aligned}
\]
\[
\boxed{\begin{aligned}
 f([-1,0])&=[0,1],\\
 f([0,2])&=[0,4],\\
 f([-1,0]\cap[0,2])&=f(0)=0,\\
 f([-1,0])\cap f([0,2])&=[0,1].
\end{aligned}}
\]

\medskip
\noindent $c)$ \underline{$\subseteq$}.
Let
\[
 x\in f^{-1}\left(\bigcup_{j\in J}Y_j\right).
\]
\noindent Then
\[
 f(x)\in\bigcup_{j\in J}Y_j
 \quad\Longrightarrow\quad
 (\exists j_0\in J):\ f(x)\in Y_{j_0}
 \quad\Longrightarrow\quad
 x\in f^{-1}(Y_{j_0})
 \subseteq\bigcup_{j\in J}f^{-1}(Y_j).
\]
\noindent \underline{$\supseteq$}.
If $x\in\bigcup_{j\in J}f^{-1}(Y_j)$, there is $j_0\in J$
with $x\in f^{-1}(Y_{j_0})$, and therefore
\[
 f(x)\in Y_{j_0}\ \Longrightarrow\ f(x)\in\bigcup_{j\in J}Y_j.
\]
% The handwritten page stops here.
.
% No preamble and no forced page break.
% EDITORIAL NOTES:
% - This page continues the work begun on Page 3: it first returns to
%   parts (a) and (b), then begins the proof of (c).
% - The equality in (b), for injective f, tacitly assumes a non-empty
%   index family; the manuscript does not discuss the empty case.
% - In the finite example A={-1,0,1}, f(x)=x^2, the handwriting seems
%   to say f(X_1)=A, although f({-1,1})={1}. The apparent A is retained.
% - Right-margin interval calculations seem to concern f(x)=x^2, but
%   the domain is not restated. The manuscript appears to say
%   f([-1,3])=[1,9]; mathematically its image would be [0,9].
%   That apparent source expression is retained and flagged here.
% - Interval brackets and certain endpoints in the boxed calculation
%   are provisional readings of the handwriting.
% - The final proof (c) ends mid-argument; no continuation or part (d)
%   has been added.

\subsection*{Seminar 2 - Page 4}

\noindent\underline{Solution (continued)}

\medskip
\noindent $a)$ \underline{$\subseteq$}.
Let $y\in f\left(\bigcup_{i\in I}X_i\right)$. Then
\[
\exists\,x\in\bigcup_{i\in I}X_i\quad\text{such that}\quad f(x)=y.
\]
\noindent Thus there is $i_0\in I$ with $x\in X_{i_0}$ and $f(x)=y$, hence
\[
 y\in f(X_{i_0})\subseteq\bigcup_{i\in I}f(X_i).
\]
\noindent \underline{$\supseteq$}.
There is $i_0\in I$ such that
$y\in f(X_{i_0})$, whence
\[
 y\in f\left(\bigcup_{i\in I}X_i\right).
\]

\medskip
\noindent $b)$ Let
\[
 y\in f\left(\bigcap_{i\in I}X_i\right).
\]
\noindent There is $x\in\bigcap_{i\in I}X_i$ with $f(x)=y$, so
\[
 y\in\bigcap_{i\in I}f(X_i).
\]
\noindent Suppose $f$ is injective. If
$y\in\bigcap_{i\in I}f(X_i)$, then
$y\in f(X_i)$ for every $i\in I$. Thus
\[
 (\forall i\in I)(\exists a_i\in X_i):\ f(a_i)=y.
\]
\noindent By injectivity, $a_i=a_j$ for all $i,j\in I$.
Denote this common element by $a$. Then
\[
 f(a)=y,\qquad a\in\bigcap_{i\in I}X_i,
 \qquad y\in f\left(\bigcap_{i\in I}X_i\right).
\]

\medskip
\noindent\underline{Example.} The manuscript also gives
\[
 A=\{-1,0,1\},\qquad B=\{0,1\},\qquad f(x)=x^2,
\]
\[
 X_1=\{-1,1\},\qquad X_2=\{-1,0\},
\]
\[
 f(X_1)=A,\qquad f(X_2)=\{0,1\}.
\]

% Small sketch of the parabola y=x^2, after the marginal drawing.
\begin{center}
\begin{tikzpicture}[>=stealth,scale=.77]
  \draw[->] (-1.8,0)--(3.0,0) node[right] {$x$};
  \draw[->] (0,-.25)--(0,3.45) node[above] {$y$};
  \foreach \x in {-1,0,1,2,3}{\draw (\x,.06)--(\x,-.06) node[below] {\scriptsize $\x$};}
  \draw[domain=-1.35:2.35,samples=70,smooth,thick] plot (\x,{.5*\x*\x});
\end{tikzpicture}
\end{center}

% Right-margin examples: interval brackets and some endpoints are uncertain.
\noindent\underline{Marginal calculations (provisional readings):}
\[
\begin{aligned}
 f([1,2])&=[1,4], & f([-1,3])&=[1,9],\\
 f([1,2])&=f([1,3]\cap[1,2])=[1,4].
\end{aligned}
\]
\[
\boxed{\begin{aligned}
 f([-1,0])&=[0,1],\\
 f([0,2])&=[0,4],\\
 f([-1,0]\cap[0,2])&=f(0)=0,\\
 f([-1,0])\cap f([0,2])&=[0,1].
\end{aligned}}
\]

\medskip
\noindent $c)$ \underline{$\subseteq$}.
Let
\[
 x\in f^{-1}\left(\bigcup_{j\in J}Y_j\right).
\]
\noindent Then
\[
 f(x)\in\bigcup_{j\in J}Y_j
 \quad\Longrightarrow\quad
 (\exists j_0\in J):\ f(x)\in Y_{j_0}
 \quad\Longrightarrow\quad
 x\in f^{-1}(Y_{j_0})
 \subseteq\bigcup_{j\in J}f^{-1}(Y_j).
\]
\noindent \underline{$\supseteq$}.
If $x\in\bigcup_{j\in J}f^{-1}(Y_j)$, there is $j_0\in J$
with $x\in f^{-1}(Y_{j_0})$, and therefore
\[
 f(x)\in Y_{j_0}\ \Longrightarrow\ f(x)\in\bigcup_{j\in J}Y_j.
\]
% The handwritten page stops here.
.
% No preamble and no forced page break.
% EDITORIAL NOTES:
% - This page continues the work begun on Page 3: it first returns to
%   parts (a) and (b), then begins the proof of (c).
% - The equality in (b), for injective f, tacitly assumes a non-empty
%   index family; the manuscript does not discuss the empty case.
% - In the finite example A={-1,0,1}, f(x)=x^2, the handwriting seems
%   to say f(X_1)=A, although f({-1,1})={1}. The apparent A is retained.
% - Right-margin interval calculations seem to concern f(x)=x^2, but
%   the domain is not restated. The manuscript appears to say
%   f([-1,3])=[1,9]; mathematically its image would be [0,9].
%   That apparent source expression is retained and flagged here.
% - Interval brackets and certain endpoints in the boxed calculation
%   are provisional readings of the handwriting.
% - The final proof (c) ends mid-argument; no continuation or part (d)
%   has been added.

\subsection*{Seminar 2 - Page 4}

\noindent\underline{Solution (continued)}

\medskip
\noindent $a)$ \underline{$\subseteq$}.
Let $y\in f\left(\bigcup_{i\in I}X_i\right)$. Then
\[
\exists\,x\in\bigcup_{i\in I}X_i\quad\text{such that}\quad f(x)=y.
\]
\noindent Thus there is $i_0\in I$ with $x\in X_{i_0}$ and $f(x)=y$, hence
\[
 y\in f(X_{i_0})\subseteq\bigcup_{i\in I}f(X_i).
\]
\noindent \underline{$\supseteq$}.
There is $i_0\in I$ such that
$y\in f(X_{i_0})$, whence
\[
 y\in f\left(\bigcup_{i\in I}X_i\right).
\]

\medskip
\noindent $b)$ Let
\[
 y\in f\left(\bigcap_{i\in I}X_i\right).
\]
\noindent There is $x\in\bigcap_{i\in I}X_i$ with $f(x)=y$, so
\[
 y\in\bigcap_{i\in I}f(X_i).
\]
\noindent Suppose $f$ is injective. If
$y\in\bigcap_{i\in I}f(X_i)$, then
$y\in f(X_i)$ for every $i\in I$. Thus
\[
 (\forall i\in I)(\exists a_i\in X_i):\ f(a_i)=y.
\]
\noindent By injectivity, $a_i=a_j$ for all $i,j\in I$.
Denote this common element by $a$. Then
\[
 f(a)=y,\qquad a\in\bigcap_{i\in I}X_i,
 \qquad y\in f\left(\bigcap_{i\in I}X_i\right).
\]

\medskip
\noindent\underline{Example.} The manuscript also gives
\[
 A=\{-1,0,1\},\qquad B=\{0,1\},\qquad f(x)=x^2,
\]
\[
 X_1=\{-1,1\},\qquad X_2=\{-1,0\},
\]
\[
 f(X_1)=A,\qquad f(X_2)=\{0,1\}.
\]

% Small sketch of the parabola y=x^2, after the marginal drawing.
\begin{center}
\begin{tikzpicture}[>=stealth,scale=.77]
  \draw[->] (-1.8,0)--(3.0,0) node[right] {$x$};
  \draw[->] (0,-.25)--(0,3.45) node[above] {$y$};
  \foreach \x in {-1,0,1,2,3}{\draw (\x,.06)--(\x,-.06) node[below] {\scriptsize $\x$};}
  \draw[domain=-1.35:2.35,samples=70,smooth,thick] plot (\x,{.5*\x*\x});
\end{tikzpicture}
\end{center}

% Right-margin examples: interval brackets and some endpoints are uncertain.
\noindent\underline{Marginal calculations (provisional readings):}
\[
\begin{aligned}
 f([1,2])&=[1,4], & f([-1,3])&=[1,9],\\
 f([1,2])&=f([1,3]\cap[1,2])=[1,4].
\end{aligned}
\]
\[
\boxed{\begin{aligned}
 f([-1,0])&=[0,1],\\
 f([0,2])&=[0,4],\\
 f([-1,0]\cap[0,2])&=f(0)=0,\\
 f([-1,0])\cap f([0,2])&=[0,1].
\end{aligned}}
\]

\medskip
\noindent $c)$ \underline{$\subseteq$}.
Let
\[
 x\in f^{-1}\left(\bigcup_{j\in J}Y_j\right).
\]
\noindent Then
\[
 f(x)\in\bigcup_{j\in J}Y_j
 \quad\Longrightarrow\quad
 (\exists j_0\in J):\ f(x)\in Y_{j_0}
 \quad\Longrightarrow\quad
 x\in f^{-1}(Y_{j_0})
 \subseteq\bigcup_{j\in J}f^{-1}(Y_j).
\]
\noindent \underline{$\supseteq$}.
If $x\in\bigcup_{j\in J}f^{-1}(Y_j)$, there is $j_0\in J$
with $x\in f^{-1}(Y_{j_0})$, and therefore
\[
 f(x)\in Y_{j_0}\ \Longrightarrow\ f(x)\in\bigcup_{j\in J}Y_j.
\]
% The handwritten page stops here.
